\bookchapter{Capitolo 61: Swift in fretta — il primo linguaggio dello sviluppo Apple}

\chapterintro{In questo capitolo impariamo davvero il linguaggio. Niente timore — Swift è famoso per essere «amichevole coi principianti» e non serve imparare a mente il manuale di sintassi: prendendo il 20\% della grammatica, leggi e scrivi l'80\% del codice quotidiano. Il resto lo cerchi mentre costruisci i progetti. Ricorda il nostro principio di studio: capire > memorizzare, saper usare > padroneggiare.}

\booksection{Variabili e tipi: dare un nome ai dati}

Un programma è «dati + logica». La variabile è un nome che dai ai dati per conservarli. Swift dichiara con let e var:

\begin{lstlisting}[language=Swift]
let name = "Marco"        // let: costante (il consigliato per default)
var age = 18             // var: variabile (si può cambiare dopo)

age = 19                 // <i class="fas fa-check-circle mr-1"></i> var si può cambiare
// name = "Anna"         // <i class="fas fa-times-circle mr-1"></i> let non si può cambiare: errore di compilazione

let score: Int = 100     // puoi anche dichiarare il tipo esplicitamente
let pi = 3.14            // se lo ometti viene dedotto automaticamente

print("Ciao, \(name)")   // interpolazione di stringa: backslash + parentesi
\end{lstlisting}

I tipi di uso quotidiano, da fissare con una tabella:

\begin{center}
\begin{tabularx}{\linewidth}{|X|X|X|}
\hline
\textbf{Tipo} & \textbf{Cosa contiene} & \textbf{Esempio} \tabularnewline
\hline
Int & Numero intero & 42 \tabularnewline
Double & Numero decimale & 3.14 \tabularnewline
String & Testo & "ciao" \tabularnewline
Bool & Vero/falso & true / false \tabularnewline
{}[Int] (array) & Un gruppo di dati & {}[1, 2, 3] \tabularnewline
\hline
\end{tabularx}
\end{center}

\begin{center}\small\color{gray}I tipi di base di uso quotidiano in Swift\end{center}

\begin{tcbinfo}{Consiglio: la regola «prima let»}
La regola ferrea della community Swift: per default usa let, e ricorri a var solo quando serve davvero. Dove basterebbe let, se metti var Xcode ti avverte in giallo. Questa abitudine ti risparmia tutta una famiglia di bug da «variabile cambiata per sbaglio» — meno cose cambiano nel codice, più è facile ragionarci.

\end{tcbinfo}

\booksection{I valori opzionali (Optional): il «talismano anti-crash» di Swift}

È il concetto più importante di Swift, e non c'è il secondo. In una frase: quando un valore può «non esserci», Swift ti obbliga a gestirlo — a livello di linguaggio fa sparire il «crash da puntatore nullo», l'incubo di tutto il mestiere.

\begin{lstlisting}[language=Swift]
var name: String = "Marco"   // tipo normale: non può mai essere nil
// name = nil                // <i class="fas fa-times-circle mr-1"></i> errore di compilazione — ti blocca sul posto

var nickname: String? = nil  // aggiungi un punto di domanda: nil consentito (opzionale)
nickname = "Nino"

// Per usare un opzionale bisogna fare l'"unwrap":
if let n = nickname {        // if let: unwrap sicuro
    print("Ciao, \(n)")
} else {
    print("Ancora nessun nickname")
}

let len = nickname?.count ?? 0   // ?? fornisce un valore predefinito (operatore di coalescenza)
print(len)
\end{lstlisting}

\begin{tcbwarn}{Attenzione: ! è un'«operazione ad alto rischio»}
Scrivere nel codice nickname! (unwrapping forzato) equivale a gridare al compilatore «ti garantisco che non è nil» — se sbagli il colpo, l'app crasha all'istante. Regola da principiante: per quanto possibile usa if let / ??, mai !. A volte Xcode, completando il codice, aggiunge un ! di sua iniziativa: ricorda di toglierlo.

\end{tcbwarn}

\booksection{Funzioni e closure: impacchettare la logica in “pulsanti”}

Una funzione è un pezzo di logica con un nome, come un pulsante sul telecomando. Le closure invece sono “funzioni senza nome”, che puoi passare a qualcuno al volo:

\begin{lstlisting}[language=Swift]
// Funzioni: prima il nome, poi parametri e tipo di ritorno
func greet(name: String) -> String {
    return "Ciao, \(name)!"
}
print(greet(name: "Marco"))     // alla chiamata serve il nome del parametro

// Closure: definite al volo, snelle e dirette
let numbers = [1, 2, 3, 4]
let doubled = numbers.map { $0 * 2 }   // $0 = elemento corrente
// [2, 4, 6, 8]

let big = numbers.filter { $0 > 2 }    // filtro
// [3, 4]
\end{lstlisting}

In SwiftUI le closure sono ovunque: le graffe dopo Button(“Toccami”) \{ ... \} sono una closure. Impara a memoria le “combo delle collezioni” map, filter e sorted: nel futuro il 90\% del lavoro con le liste sarà una loro combinazione.

\booksection{Struct e protocolli: il sistema “LEGO” di Swift}

Le preferenze ufficiali di Swift sono chiare: prima gli struct, e le class solo quando serve l'ereditarietà. Gli struct sono “tipi valore”: quando li copi, ognuno vive di vita propria, eliminando di natura tutta una categoria di bug:

\begin{lstlisting}[language=Swift]
// La struct: una scatola per dati
struct Record {
    var id = UUID()
    var amount: Double
    var category: String
    var date: Date

    // Le struct hanno un "memberwise initializer" integrato: si creano subito
}
let r = Record(amount: 25.5, category: "Pasti", date: Date())

// Il protocollo: un patto su "quali capacità sono richieste"
protocol Identifiable {
    var id: UUID { get }
}
struct Record: Identifiable { ... }   // conformarsi a un protocollo = promettere di soddisfarne i requisiti
\end{lstlisting}

In SwiftUI quasi tutto è uno struct (View è un protocollo). Se per ora non ti è chiaro, nessun problema: dal capitolo 62 ci ritroverai ogni giorno. Nel mondo di Swift, abbraccia subito gli struct.

\booksection{La scheda rapida di sintassi per chi parte da zero}

\begin{center}
\begin{tabularx}{\linewidth}{|X|X|}
\hline
\textbf{Cosa vuoi fare} & \textbf{Come si scrive} \tabularnewline
\hline
Definire una costante/variabile & \texttt{let x = 1}/\texttt{var x = 1} \tabularnewline
Interpolazione di stringhe & "Ciao, \textbackslash{}(name)" \tabularnewline
Definire una funzione & func nome(parametro: Tipo) -> TipoRisultato \{ ... \} \tabularnewline
Prendere decisioni & if condizione \{ ... \} else \{ ... \} \tabularnewline
Iterare & for x in array \{ ... \} \tabularnewline
Unwrapping sicuro & if let x = valoreOpzionale \{ ... \} \tabularnewline
Valore predefinito & x ?? valorePredefinito \tabularnewline
Definire uno struct & \texttt{struct X \{ var a: Int; var b: String \}} \tabularnewline
Closure & \texttt{\{ \$0 * 2 \}} \tabularnewline
\hline
\end{tabularx}
\end{center}

\begin{center}\small\color{gray}Scheda rapida della sintassi Swift (versione principianti)\end{center}

\begin{tcbpractice}{Mettiti alla prova: gioca 15 minuti con Swift nei Playgrounds}
Con un Mac apri Xcode Playground; senza Mac usa Swift Playgrounds su iPad oppure l'arena di prova sul web. Esegui ogni blocco di codice di questo capitolo, cambia numeri e testi, provoca di proposito qualche errore e guarda come te lo segnala Xcode. Questi 15 minuti di “prima giocare, poi imparare” ti renderanno il capitolo 62 due volte più facile.

\end{tcbpractice}

\begin{tcbsummary}{Riepilogo del capitolo}
\begin{itemize}
  \item Prima let, var con cautela; i tipi si deducono da soli, ma puoi anche dichiararli esplicitamente dopo i due punti
  \item Gli optional sono l'anima di Swift: ciò che può essere vuoto va gestito; if let e ?? sono il tuo amuleto, ! è un'operazione ad alto rischio
  \item Nelle chiamate di funzione va indicato il nome dei parametri; closure = funzione senza nome, presente ovunque in SwiftUI
  \item Prima gli struct, conformandosi ai protocolli; in SwiftUI quasi tutto è uno struct
  \item Il 20\% della sintassi fa l'80\% del lavoro: prima impara a usarla, poi perfezionala
\end{itemize}

\end{tcbsummary}
